\documentclass[11pt]{article}
\usepackage[margin=25mm]{geometry}
\usepackage{fontspec}
\setmainfont{Latin Modern Roman}
\usepackage{amsmath,amssymb,mathtools}
\usepackage{unicode-math}
\setmathfont{Latin Modern Math}
\usepackage{xurl,hyperref}
\hypersetup{colorlinks=true,urlcolor=blue}
\setcounter{secnumdepth}{0}
\setlength{\parindent}{0pt}
\setlength{\parskip}{5pt}
\setlength{\emergencystretch}{3em}
\newcommand{\sref}[2]{\hyperlink{source:#1:#2}{[S#2]}}
\newcommand{\cataloguescope}[1]{\paragraph{Recorded target.}#1}
\begin{document}
\section{Erdős-Hajnal Conjecture}\label{statement}

\subsection{Definitions and mathematical statement}

\paragraph{Shared notation.}$\mathbb N=\{1,2,\ldots\}$, and $\mathbb Z,\mathbb Q,\mathbb R,\mathbb C$ denote the integers, rationals, reals and complex numbers. Any other conventions are specified locally.


All graphs are finite, simple and undirected. For $X\subseteq V(G)$ the induced subgraph $G[X]$ has exactly the edges of $G$ whose two endpoints lie in $X$. For a fixed graph $H$, say $G$ is $H$-free if no induced subgraph of $G$ is isomorphic to $H$. A clique is a set of pairwise adjacent vertices, and an independent set is a set of pairwise nonadjacent vertices. Their largest cardinalities are $\omega(G)$ and $\alpha(G)$, respectively.
\paragraph{Statement recorded in the supplied source.}
For every fixed finite graph $H$, does there exist a real number $\varepsilon(H)>0$ such that, for every $n\ge1$ and every $n$-vertex $H$-free graph $G$,
\[
 \max\{\omega(G),\alpha(G)\}\ge n^{\varepsilon(H)}?
\]
The exponent may depend on $H$ but must be independent of $n$ and $G$. Here forbidden subgraphs are induced, not arbitrary subgraphs. \sref{statement}{2}
\subsection{Short English statement}
For each fixed forbidden induced graph, must every graph avoiding it contain a clique or independent set whose size grows as a positive power of the number of vertices?
\subsection{Sources}
\begin{description}
\item[\hypertarget{source:statement:1}{S1}] ULAM.AI and the UnsolvedMath Contributors, UnsolvedMath: A Curated Collection of Open Mathematics Problems, catalogue entry. CC BY 4.0. \url{https://huggingface.co/datasets/ulamai/UnsolvedMath}.
\item[\hypertarget{source:statement:2}{S2}] Shenwei Huang, Yiao Ju and Yidong Zhou, Erdős–Hajnal Conjecture Beyond Five-Vertex Graphs, arXiv:2606.06258v3 (2026), Section 1, Conjecture 1, pages 1--2. \url{https://arxiv.org/pdf/2606.06258v3}.
\end{description}
\label{end:statement}
\end{document}
